% ============================================================================
% Fichier  : partie4/P04_C14_audit_application.tex
% Titre    : Auditer une Application en Service
% Auteur   : MINKA MI NGUIDJOÏ Thierry Emmanuel
% Version  : 1.0 -- novembre 2026
% Base   : ACSI C14, repris et développé pour INF3226
% ============================================================================
% Sources : guide d'audit des SI du CHAI (2014), fiche 3.4 « Audit des
% applications informatiques en service » (deux visées : fiabilité et
% sécurité / efficacité et performance ; points de contrôle sur les
% entrées, rejets, traitements, sorties, interfaces, données sensibles,
% piste d'audit), paraphrasée ; TP « application » du cours Audit_SI de
% l'auteur (tests SQL : doublons, équilibre des journaux, seuil de
% validation), repris et mis en cohérence (OUI/NON/NEK, observation en
% sept éléments). Faits TRANSLOG SERVICES repris du chapitre 4.
% ============================================================================

\chapter{Auditer une Application en Service}
\label{chap:audit-applications}
\chaptermark{Auditer une application}

\epigraph{%
    « Le contrôle programmé dit ce que l'application devrait faire.\\
    Les données disent ce qu'elle a fait. »%
}{Adage des auditeurs de données}

\niveauChapitre{Intermédiaire}{%
    Maîtrise du Chapitre~\ref{chap:controle-interne} (contrôles
    d'application), du Chapitre~\ref{chap:preuve} (preuves et extractions)
    et du Chapitre~\ref{chap:conclure} (observation)%
}

\begin{boxobjectifs}
\begin{itemize}
    \item Distinguer l'audit de fiabilité et l'audit d'efficacité d'une
          application.
    \item Décomposer l'audit d'une application selon le modèle
          entrées-traitements-sorties, complété des interfaces et des
          données permanentes.
    \item Tester un contrôle d'application, après avoir vérifié les
          contrôles généraux qui le protègent.
    \item Réaliser des tests substantifs sur les données~: doublons,
          équilibre des journaux, seuils, séquences, rapprochements.
    \item Rédiger l'observation d'un audit d'application.
\end{itemize}
\end{boxobjectifs}

\bigskip

% ============================================================================
\section*{Question directrice : l'application fait-elle ce qu'on croit~?}
% ============================================================================

L'application intégrée de LiZbiZ gère les expéditions, l'entreposage, la
facturation et la comptabilité. Le DAF, Mme~MBALLA, assure que « le
système bloque tout »~: les doublons, les montants au-dessus de
5~millions FCFA sans double validation, les factures sans expédition. Le
Chapitre~\ref{chap:fraude} a pourtant montré 23~paiements de
4,9~millions FCFA à un même fournisseur.

\medskip
\textbf{Plan du chapitre.}
\begin{enumerate}
    \item Décomposer l'audit d'une application (section~\ref{sec:c18-decomposer}).
    \item Tester les contrôles d'application (section~\ref{sec:c18-controles}).
    \item Les contrôles des entrées, en détail (section~\ref{sec:c18-entrees}).
    \item Interroger les données (section~\ref{sec:c18-donnees}).
    \item Analyser les données en Python (section~\ref{sec:c18-python}).
    \item Rédiger les observations (section~\ref{sec:c18-observations}).
\end{enumerate}

\separateur

% ============================================================================
\section{Décomposer l'Audit d'une Application}
\label{sec:c18-decomposer}
% ============================================================================

\subsection{Deux visées}

Le guide du CHAI distingue deux visées~\cite{chai2014guide}~:
\begin{itemize}
    \item l'audit de \textbf{fiabilité et de sécurité}, qui apprécie le
          contrôle interne de l'application et la validité des données
          qu'elle traite et restitue~;
    \item l'audit d'\textbf{efficacité et de performance}, qui apprécie son
          adéquation aux besoins, sa contribution à la valeur, sa
          rentabilité et sa pérennité.
\end{itemize}
Ce chapitre traite surtout de la première. La seconde rejoint l'audit des
projets (Chapitre~\ref{chap:audit-projets}). Un signal les relie
cependant~: lorsque les utilisateurs tiennent des tableurs en parallèle de
l'application, c'est souvent qu'elle ne répond pas à leurs besoins, et ces
tableurs échappent à tous ses contrôles.

\subsection{Entrées, traitements, sorties, interfaces, paramètres}

Une application reçoit des données, les traite et produit des résultats.
Chaque étape porte ses risques et ses contrôles. Deux éléments complètent
ce modèle classique~: les \emph{interfaces}, par où les données passent
d'un module ou d'une application à l'autre, et les \emph{données
permanentes et paramètres} (fiches fournisseurs, barèmes, seuils), dont
une seule modification change le traitement de milliers d'opérations.

\begin{tableaucours}{cL{2.8cm}L{4.4cm}L{4.0cm}}{%
    \tableauHeader{Code} &
    \tableauHeader{Sous-domaine} &
    \tableauHeader{Risques} &
    \tableauHeader{Contrôles attendus}
}
AP.1 & Entrées & Saisie erronée, incomplète, en double, non autorisée &
Contrôles de format et de vraisemblance~; détection des doublons~;
autorisation~; rejets conservés et recyclés \\
AP.2 & Traitements & Calcul faux, traitement incomplet, écriture
déséquilibrée & Équilibre automatique des écritures~; totaux de contrôle
de fin de journée~; comptes rendus des traitements par lots \\
AP.3 & Sorties & État incomplet, diffusé à la mauvaise personne & Séquence
des documents~; validation et diffusion contrôlées des états \\
AP.4 & Interfaces & Données perdues ou altérées entre modules &
Rapprochement des volumes et des montants entre modules~; mêmes contrôles
pour les fichiers que pour la saisie \\
AP.5 & Données permanentes et paramètres & Fournisseur fictif, seuil
modifié & Création des tiers sous double contrôle~; piste d'audit des
paramètres (auteur, date, valeur avant et après) \\
\end{tableaucours}
\vspace{-0.8em}
\begin{center}{\small\itshape Tableau~18.1. Décomposition de l'audit du
cycle achats-paiements (d'après la fiche 3.4 du guide du CHAI)}\end{center}

Les habilitations, sixième volet du guide du CHAI pour une application,
relèvent du domaine~D (Chapitre~\ref{chap:decomposition}) et de la matrice
des droits (Chapitre~\ref{chap:controle-interne}, ACSI-LAB-02)~: on ne les
recompte pas ici, mais on en tire les conséquences.

% ============================================================================
\section{Tester les Contrôles d'Application}
\label{sec:c18-controles}
% ============================================================================

\subsection{Les contrôles généraux d'abord}

Le Chapitre~\ref{chap:controle-interne} l'a établi~: un contrôle
d'application ne vaut que si les contrôles généraux qui le protègent sont
efficaces. Avant de tester le blocage des paiements de plus de 5~millions,
l'auditeur vérifie qui peut modifier ce seuil, et si la modification
laisse une trace. Chez LiZbiZ, le seuil est stocké dans une table de
paramètres que M.~KOUNOU peut modifier directement, et cette table n'a
aucune piste d'audit.

\subsection{Trois façons de tester un contrôle programmé}

\begin{tableaucours}{L{2.8cm}L{5.2cm}L{4.2cm}}{%
    \tableauHeader{Test} &
    \tableauHeader{Principe} &
    \tableauHeader{Exemple chez LiZbiZ}
}
Inspection du paramétrage (DOC) & Lire la configuration du contrôle dans
l'application & Le seuil de double validation vaut 5~000~000 FCFA \\
Ré-exécution (TEST) & Saisir une transaction d'essai, en environnement de
recette, qui doit être rejetée & Un paiement de 5,1~millions est bloqué
sans seconde signature \\
Test sur les données (TEST) & Chercher dans la production les opérations
que le contrôle aurait dû empêcher & Paiements validés par une seule
personne au-dessus du seuil~; paiements juste sous le seuil \\
\end{tableaucours}
\vspace{-0.8em}
\begin{center}{\small\itshape Tableau~18.2. Tester un contrôle programmé}\end{center}

Les deux premiers tests établissent que le contrôle est \emph{conçu} et
\emph{en place} aujourd'hui. Seul le troisième dit s'il a \emph{fonctionné}
sur la période, puisqu'un paramètre a pu être modifié puis rétabli. C'est
pourquoi le test sur les données est le plus probant.

\begin{boxalert}
\textbf{Violation typique.} Ré-exécuter un contrôle en production, avec une
vraie transaction, « pour voir ».

\smallskip
\textbf{Règle~:} la ré-exécution se fait en environnement de recette, sur
autorisation écrite, avec une configuration identique à la production
(vérifiée). En production, l'auditeur lit et extrait~; il ne saisit pas.
\end{boxalert}

% ============================================================================
\section{Les Contrôles des Entrées, en Détail}
\label{sec:c18-entrees}
% ============================================================================

La plupart des anomalies naissent à la saisie. Les contrôles des entrées
se classent en six types, que l'auditeur reconnaît dans le paramétrage.

\begin{tableaucours}{L{3.0cm}L{4.6cm}L{4.6cm}}{%
    \tableauHeader{Type} &
    \tableauHeader{Principe} &
    \tableauHeader{Exemple (achats de LiZbiZ)}
}
Format & Le champ respecte une syntaxe & Le RIB a la longueur et la clé
attendues \\
Plage & La valeur est dans des bornes & Quantité commandée entre 1 et
10~000 \\
Existence & La valeur renvoie à une donnée de référence & Le fournisseur
existe et n'est pas bloqué \\
Cohérence & Deux champs sont compatibles & Date de facture postérieure à
la date de commande \\
Unicité & Pas de doublon & Couple fournisseur et numéro de facture unique
(après normalisation) \\
Total de contrôle & Un total calculé est comparé à un total saisi &
Total d'un lot de factures comparé au bordereau \\
\end{tableaucours}
\vspace{-0.8em}
\begin{center}{\small\itshape Tableau~18.3. Les six types de contrôles
des entrées}\end{center}

\begin{boxalert}
\textbf{Violation typique.} Constater que l'application refuse un doublon
exact de numéro de facture, et conclure que les doublons sont impossibles.

\smallskip
\textbf{Règle~:} un contrôle d'unicité ne vaut que sur des données
normalisées. « FA-102 », « FA102 » et « fa 102 » sont trois valeurs pour
l'application, une seule facture pour le fournisseur~: les six doublons
de LiZbiZ (31,4~millions FCFA) en sont la preuve.
\end{boxalert}

% ============================================================================
\section{Interroger les Données}
\label{sec:c18-donnees}
% ============================================================================

\subsection{Les techniques d'audit assistées par ordinateur}

Les tests substantifs sur les données portent sur la population entière,
et non sur un échantillon~: ils suppriment le risque d'échantillonnage
(Chapitre~\ref{chap:preuve}). Ils reposent sur une extraction tracée, avec
lettre et empreinte. Le Chapitre~\ref{chap:audit-continu} présentera les
outils~; voici les cinq tests de base, en SQL, sur les tables
\texttt{PAIEMENT}, \texttt{ECRITURE} et \texttt{EXPEDITION} d'une copie de
la base.

\paragraph{Doublons.} Un même fournisseur, un même montant, une même
facture, payés deux fois.
\begin{lstlisting}[language=SQL]
SELECT id_fournisseur, montant, ref_facture, COUNT(*) AS nb
FROM PAIEMENT
GROUP BY id_fournisseur, montant, ref_facture
HAVING COUNT(*) > 1;
\end{lstlisting}

Le contrôle de l'application compare souvent la référence exacte. Une
variante («~FA-102~» puis «~FA102~») passe~: l'auditeur normalise donc les
références (majuscules, sans espaces ni tirets) avant de grouper.

\paragraph{Équilibre des journaux.} Total des débits égal au total des
crédits, par pièce et par journal.
\begin{lstlisting}[language=SQL]
SELECT code_journal, num_piece, SUM(debit) AS d, SUM(credit) AS c
FROM ECRITURE
GROUP BY code_journal, num_piece
HAVING SUM(debit) <> SUM(credit);
\end{lstlisting}

\paragraph{Seuils.} Opérations au-dessus du seuil validées par une seule
personne, et opérations concentrées juste sous le seuil.
\begin{lstlisting}[language=SQL]
SELECT id_fournisseur, COUNT(*) AS nb, SUM(montant) AS total
FROM PAIEMENT
WHERE montant >= 4500000 AND montant < 5000000
GROUP BY id_fournisseur
ORDER BY nb DESC;
\end{lstlisting}

\paragraph{Séquences.} Numéros de factures manquants ou en double dans la
séquence.

\paragraph{Rapprochements entre modules.} Chaque expédition livrée a-t-elle
une facture~? C'est le test de l'interface entre l'expédition et la
facturation.
\begin{lstlisting}[language=SQL]
SELECT e.num_expedition, e.date_livraison
FROM EXPEDITION e
LEFT JOIN FACTURE f ON f.num_expedition = e.num_expedition
WHERE e.statut = 'LIVREE' AND f.num_facture IS NULL;
\end{lstlisting}

\subsection{Ce que disent les données de LiZbiZ}

\begin{boxscenario}{Opération LIZBIZ. Les cinq tests de 2025}
\textbf{Population.} 4~812~paiements fournisseurs, 61~240~lignes
d'écritures, 18~340~expéditions livrées (extraction du 15~mai~2026, empreinte
au papier de travail PT-21).

\medskip
\textbf{Résultats.}
\begin{itemize}
    \item \emph{Doublons}~: 6~paiements payés deux fois, pour 31,4~millions
          FCFA, tous avec une référence de facture écrite de deux façons.
    \item \emph{Équilibre}~: aucune pièce déséquilibrée.
    \item \emph{Seuils}~: 41~paiements entre 4,5 et 5~millions, dont 23 au
          seul fournisseur TRANSLOG SERVICES (Chapitre~\ref{chap:fraude})~;
          aucun paiement au-dessus du seuil sans seconde signature.
    \item \emph{Séquences}~: la numérotation des factures est continue.
    \item \emph{Rapprochement}~: 14~expéditions livrées sans facture, pour un
          chiffre d'affaires estimé à 9,8~millions FCFA.
\end{itemize}

\medskip
\textbf{Analyse} (confiance élevée sur les faits, moyenne sur les causes).
Le contrôle du seuil a fonctionné au sens strict~: aucun paiement ne l'a
dépassé sans seconde signature. Il a été contourné par fractionnement.
Le contrôle des doublons existe, mais il est trop étroit. L'interface
expéditions-facturation perd des opérations.
\end{boxscenario}

\subsection{Les tests lèvent l'incertitude}

\begin{boxlecture}{Cycle achats-paiements}
Avant les tests sur les données, AUDIX CONSEIL a coté les vingt contrôles
du cycle (AP.1 à AP.5) par inspection et entretiens~: 12~\OUI{}, 3~\NON{},
5~\NEK{}. Les cinq \NEK{} portaient sur ce que seules les données pouvaient
établir~: la détection effective des doublons, l'équilibre des écritures,
la continuité des séquences, l'exhaustivité de l'interface
expéditions-facturation et la piste d'audit des paramètres.

\medskip
\begin{tableaucours}{L{4.0cm}cccccc}{%
    \tableauHeader{Moment} &
    \tableauHeader{\OUI{}} & \tableauHeader{\NON{}} & \tableauHeader{\NEK{}} &
    \tableauHeader{$\mu$} & \tableauHeader{$U$} & \tableauHeader{$H$}
}
Inspection et entretiens & 12 & 3 & 5 & 0,60 & 0,25 & 1,35 \\
Après les tests sur données & 14 & 6 & 0 & 0,70 & 0,00 & 0,88 \\
\end{tableaucours}

\medskip
\textbf{Constatation.} Deux \NEK{} sont levés en \OUI{} (équilibre,
séquences) et trois en \NON{} (doublons, interface, piste d'audit des
paramètres).

\medskip
\textbf{Analyse} (confiance élevée). Les tests sur données ont fait passer
le cycle d'une conformité opaque à une conformité claire, avec six
non-conformités connues. Le taux de conformité a augmenté \emph{et} le
nombre de \NON{} a doublé~: c'est ce que produit la levée de
l'incertitude.

\medskip
\hfill\lecture{0{,}70}{0{,}00}{0{,}88}
\end{boxlecture}

% ============================================================================
\section{Analyser les Données en Python}
\label{sec:c18-python}
% ============================================================================

Le langage SQL suffit pour la plupart des tests. Python devient utile
lorsqu'il faut normaliser des valeurs, enchaîner plusieurs fichiers ou
produire un graphique. La bibliothèque \texttt{pandas} manipule des
tableaux comme une base de données.

\begin{boxoutils}{Doublons après normalisation, en Python}
\begin{lstlisting}[language=Python]
import pandas as pd, re
p = pd.read_csv("paiements_fournisseurs_2025.csv", sep=";")
# normaliser le numéro de facture : majuscules, sans espace ni tiret
p["ref_norm"] = p["ref_facture"].str.upper().map(
    lambda s: re.sub(r"[^A-Z0-9]", "", s))
cle = ["id_fournisseur", "ref_norm", "montant_fcfa"]
doublons = p[p.duplicated(cle, keep=False)].sort_values(cle)
print(len(doublons) // 2, doublons["montant_fcfa"].sum() // 2)
\end{lstlisting}
\end{boxoutils}

\begin{boxoutils}{Ruptures de séquence, en SQL}
\begin{lstlisting}[language=SQL]
-- numéros de bons de réception manquants dans la séquence
SELECT numero + 1 AS debut_trou,
       suivant - 1 AS fin_trou
FROM (SELECT numero,
             LEAD(numero) OVER (ORDER BY numero) AS suivant
      FROM   receptions)
WHERE suivant - numero > 1;
\end{lstlisting}
\end{boxoutils}

\begin{boxprocedure}
\textbf{Rendre une analyse reproductible}
\begin{enumerate}
    \item Calculer l'empreinte des fichiers reçus avant toute ouverture
          (Chapitre~\ref{chap:preuve}).
    \item Écrire l'analyse dans un script, et non par manipulations à la
          main dans un tableur.
    \item Conserver le script, sa version, et sa sortie dans les papiers de
          travail.
    \item Vérifier les totaux de contrôle~: nombre de lignes et montant total
          lus doivent égaler ceux de la lettre d'extraction.
\end{enumerate}
\end{boxprocedure}

% ============================================================================
\section{Rédiger les Observations}
\label{sec:c18-observations}
% ============================================================================

Une observation d'audit d'application relie un fait tiré des données à un
contrôle défaillant, et la recommandation porte sur le contrôle, non sur
les seules opérations découvertes. Récupérer les 31,4~millions ne suffit
pas~: il faut que le prochain doublon soit bloqué.

\begin{boxobjet}{AP.1.3~: Observation}
\textbf{1. Norme~:} procédure des paiements fournisseurs de LiZbiZ~:
« une facture ne peut être payée qu'une fois »~; contrôle programmé de
détection des doublons.

\textbf{2. Constatation~:} sur 4~812~paiements de 2025, 6 factures ont été
payées deux fois, pour 31,4~millions FCFA. Dans les six cas, la référence
de la facture est saisie de deux façons («~FA-102~» et «~FA102~»).
Test sur la population entière, papier de travail PT-22. Cote \NON{}.

\textbf{3. Risques~:} paiement indu~; perte financière~; usage du défaut
pour détourner des fonds.

\textbf{4. Contradictoire~:} le DAF indique, le 3~juin~2026, que quatre des
six doublons ont été récupérés par avoir en 2025~; pièces remises.

\textbf{5. Analyse} (confiance élevée)~: le contrôle compare la référence
exacte. Il ne normalise pas la saisie et ne compare pas le fournisseur, le
montant et la date. La récupération de quatre doublons montre que le
défaut était connu de la comptabilité, sans être remonté.

\textbf{6. Conclusion~:} le contrôle de détection des doublons n'est pas
efficace. Deux paiements, 10,6~millions FCFA, restent à récupérer.

\textbf{7. Recommandations~:} normaliser les références avant comparaison
et bloquer tout paiement de même fournisseur et même montant dans un
délai de 60~jours sans validation spécifique~; revue mensuelle des
doublons potentiels par le contrôle interne~; récupération des deux
paiements restants.
\end{boxobjet}

\begin{boxmetier}{Contrôleur interne}
\textbf{Installer les tests sur données dans le contrôle permanent}
\begin{enumerate}
    \item Obtenir un accès en lecture seule à une copie de la base, mise à
          jour chaque mois.
    \item Programmer les cinq tests de base~: doublons, équilibre, seuils,
          séquences, rapprochements.
    \item Fixer pour chaque test un seuil d'alerte et un destinataire.
    \item Documenter chaque anomalie~: explication, correction, ou
          escalade.
    \item Remettre les résultats à l'auditeur~: ils deviennent des preuves
          de fonctionnement du contrôle permanent.
\end{enumerate}
\end{boxmetier}

\begin{boxinfo}
\textbf{Le commissaire aux comptes.}
\textbf{Tirer parti des tests de l'auditeur SI}
\begin{enumerate}
    \item Vérifier l'extraction~: source, date, empreinte, exhaustivité
          rapprochée des totaux comptables.
    \item Reprendre les doublons non récupérés et les expéditions non
          facturées dans l'évaluation des anomalies.
    \item Apprécier si le fractionnement des paiements de TRANSLOG SERVICES
          appelle les diligences du Chapitre~\ref{chap:droit}.
\end{enumerate}
\end{boxinfo}

% ============================================================================

\separateur

\begin{boxresume}
\textbf{Ce qu'il faut retenir de ce chapitre~:}
\begin{itemize}
    \item \textbf{Deux visées}~: fiabilité et sécurité, efficacité et
          performance~; un tableur parallèle est un signal pour les deux.
    \item \textbf{Cinq sous-domaines}~: entrées, traitements, sorties,
          interfaces, données permanentes et paramètres~; les habilitations
          relèvent du domaine~D.
    \item \textbf{Les contrôles généraux d'abord}~: un seuil modifiable sans
          trace ne protège rien.
    \item \textbf{Inspection et ré-exécution} montrent qu'un contrôle est en
          place~; seul le test sur les données dit qu'il a fonctionné.
    \item \textbf{Cinq tests de base}~: doublons, équilibre, seuils,
          séquences, rapprochements entre modules.
    \item \textbf{Les tests sur données lèvent l'incertitude}~: le taux de
          conformité peut monter alors que les \NON{} augmentent.
    \item \textbf{Les contrôles des entrées} (format, plage, existence, cohérence, unicité, total) se testent sur des données normalisées~; une analyse de données est un script conservé, pas une manipulation de tableur.
\end{itemize}
\end{boxresume}

\bigskip

\begin{boxexo}[title={\ding{45}\quad Exercice~18.1. Rattacher les contrôles \niveaubase}]
Rattachez chaque contrôle à un sous-domaine (AP.1 à AP.5).
\begin{enumerate}[(a)]
    \item Les factures rejetées sont conservées et recyclées sous cinq jours.
    \item Toute création de fournisseur est validée par une seconde personne.
    \item Les totaux de fin de journée sont rapprochés du journal.
    \item Les états de paiement ne sont transmis qu'à la trésorerie.
    \item Le nombre de factures transmises par le module de facturation est
          rapproché du nombre reçu par la comptabilité.
\end{enumerate}
\end{boxexo}

\begin{boxexo}[title={\ding{45}\quad Exercice~18.2. Écrire les requêtes \niveaumoyen}]
Sur les tables \texttt{PAIEMENT} et \texttt{FOURNISSEUR} de LiZbiZ~:
\begin{enumerate}[(a)]
    \item Écrivez la requête qui trouve les paiements au-dessus de
          5~millions validés par la même personne que celle qui les a
          saisis.
    \item Écrivez la requête qui trouve les fournisseurs dont les
          coordonnées bancaires sont identiques à celles d'un salarié.
    \item Pourquoi faut-il normaliser les références de facture avant de
          chercher les doublons~?
    \item Quelle cotation retenir pour le contrôle des doublons si
          l'extraction ne couvre que neuf mois de l'exercice~?
\end{enumerate}
\end{boxexo}

\begin{boxexo}[title={\ding{45}\quad Exercice~18.3. L'interface expéditions-facturation \niveauexpert}]
Reprenez le résultat~: 14~expéditions livrées sans facture, 9,8~millions
FCFA.
\begin{enumerate}[(a)]
    \item Rédigez l'observation en sept éléments.
    \item Quelles hypothèses sur la cause faut-il vérifier (défaut technique,
          expéditions gratuites non documentées, fraude), et avec quelles
          preuves~?
    \item Proposez un contrôle permanent de l'interface, et sa fiche
          d'assertion.
    \item Quel indicateur suivrait l'efficacité de ce contrôle~?
\end{enumerate}
\end{boxexo}

\bigskip

\begin{boxexo}[title={\ding{45}\quad Exercice~18.4. Les types de contrôles \niveaubase}]
Classez selon le Tableau~18.3~: (a) refus d'une facture datée de 2027~;
(b) refus d'un fournisseur inconnu~; (c) refus d'un RIB de 20~caractères au
lieu de 23~; (d) comparaison du total d'un lot avec son bordereau.
\end{boxexo}

\begin{boxexo}[title={\ding{45}\quad Exercice~18.5. Une requête à écrire \niveaumoyen}]
Écrivez la requête qui liste les factures payées sans bon de réception
correspondant (tables \texttt{factures}, \texttt{receptions},
\texttt{paiements}). Quels faux positifs attendez-vous (prestations de
service, acomptes)~?
\end{boxexo}

\begin{boxexo}[title={\ding{45}\quad Exercice~18.6. Normaliser \niveaumoyen}]
Expliquez pourquoi la normalisation du script Python retrouve des doublons
que l'application ne voit pas. Quel risque inverse crée-t-elle (deux
factures différentes rendues identiques)~? Comment l'écarter~?
\end{boxexo}

\begin{boxexo}[title={\ding{45}\quad Exercice~18.7. La séquence \niveaumoyen}]
La requête des ruptures de séquence trouve un trou de 12 numéros en
décembre. Que faites-vous avant de conclure~? Quelles explications
légitimes peut-il y avoir~?
\end{boxexo}

\begin{boxexo}[title={\ding{45}\quad Exercice~18.8. Les paramètres \niveaumoyen}]
Le seuil de double validation de 5~millions est stocké dans une table de
paramètres. Rédigez la fiche d'assertion du contrôle « le seuil n'est
modifié que par une procédure autorisée » et le test associé.
\end{boxexo}

\begin{boxexo}[title={\ding{45}\quad Exercice~18.9. Reproductibilité \niveauexpert}]
Un auditeur a trouvé les 41~paiements entre 4,5 et 5~millions par un
filtre dans un tableur, sans conserver le fichier filtré. Le contradictoire
conteste le chiffre. Que fallait-il conserver, et comment refaites-vous le
test de façon opposable~?
\end{boxexo}

\begin{boxexo}[title={\ding{45}\quad Exercice~18.10. Les factures sans expédition \niveauexpert}]
Sur 18~377~expéditions, 37 sont annulées et 14 livrées sans facture, pour
9,8~millions FCFA. Rédigez l'observation en sept éléments et proposez un
contrôle d'application qui empêcherait le cas.
\end{boxexo}

\bigskip

\begin{boxinfo}
\textbf{Points de contrôle de l'élève.} Je sais~: décomposer une
application en entrées, traitements, sorties, interfaces et paramètres~;
tester un contrôle programmé~; classer un contrôle des entrées~; écrire les
requêtes de doublons, de seuils et de séquences~; rendre une analyse
reproductible.
\end{boxinfo}

\begin{boxinfo}
\textbf{Pour aller plus loin.}
\begin{itemize}
    \item INTOSAI, GUID~5100~\cite{intosai2019guid5100}~: les techniques
          d'audit assistées par ordinateur.
    \item Le guide du CHAI~\cite{chai2014guide}~: les fiches « applications
          en service ».
    \item M.~Nigrini~\cite{nigrini2012benford}~: l'analyse de données
          pour l'audit.
\end{itemize}
\end{boxinfo}

\begin{boxinfo}
\textbf{Transition.} Une application en service a d'abord été un projet.
Le Chapitre~\ref{chap:audit-projets} audite les projets informatiques~:
enjeux, pilotage, tests, mise en production, et les signaux d'un projet en
dérive.
\end{boxinfo}
